\documentclass[10pt]{article}

\usepackage[T1]{fontenc}
\usepackage[utf8]{inputenc}
\usepackage{lmodern}
\usepackage[letterpaper,margin=0.72in]{geometry}
\usepackage{amsmath,amssymb}
\usepackage{booktabs,tabularx,array,longtable}
\usepackage{float}
\usepackage{xcolor}
\usepackage{microtype}
\usepackage{enumitem}
\usepackage{fancyhdr}
\usepackage{hyperref}
\usepackage{graphicx}

\definecolor{navy}{HTML}{163A5F}
\definecolor{blue}{HTML}{2C6EAA}
\definecolor{lightblue}{HTML}{EAF3FA}
\definecolor{lightgray}{HTML}{F3F5F7}
\definecolor{green}{HTML}{13795B}
\definecolor{orange}{HTML}{B05B13}

\hypersetup{
  colorlinks=true,
  linkcolor=navy,
  citecolor=blue,
  urlcolor=blue,
  pdftitle={RelaySpec: End-to-End Research Report},
  pdfauthor={Anonymous}
}

\pagestyle{fancy}
\fancyhf{}
\lhead{\textcolor{navy}{RelaySpec}}
\rhead{\textcolor{navy}{Research report -- August 2026}}
\cfoot{\thepage}
\setlength{\headheight}{14pt}

\setlist[itemize]{leftmargin=1.25em,itemsep=2pt,topsep=3pt}
\setlist[enumerate]{leftmargin=1.4em,itemsep=2pt,topsep=3pt}
\renewcommand{\arraystretch}{1.12}

\newcommand{\method}{\textsc{RelaySpec}}
\newcommand{\sourcebase}{\textsc{Source-Reuse}}
\newcommand{\good}[1]{\textbf{\textcolor{green}{#1}}}
\newcommand{\code}[1]{\texttt{\detokenize{#1}}}
\newcommand{\callout}[2]{%
  \begin{center}
  \colorbox{lightblue}{\parbox{0.94\linewidth}{\textbf{\textcolor{navy}{#1}}\\[-2pt]#2}}
  \end{center}
}

\title{\vspace{-1.2em}\textbf{\textcolor{navy}{RelaySpec}}\\
\large Reusing Frozen Feature-Conditioned Proposers Across Target Scales\\[3pt]
\normalsize End-to-End Method, Evidence, Amdahl Analysis, and ICLR Paper Plan}
\author{Anonymous research report}
\date{August 28, 2026}

\begin{document}
\maketitle
\vspace{-1.4em}

\begin{abstract}
Feature-conditioned speculative drafters such as DFlash can accelerate a
target language model, but their conditioning interface is normally tied to
the target used during draft training. A naive cross-scale reuse path can run
the original small source transformer on every speculative cycle to reproduce
that interface; our profiles show that this redundant trunk consumes
27--40\% of end-to-end request time. \method{} replaces the repeated source
trunk with a single learned translator from hidden features already computed
by the exact target verifier into the frozen proposer's existing conditioning
space. The target, source model, and proposer remain frozen, and ordinary
speculative verification still controls every committed token. On Qwen3-8B
and Qwen3-14B with the released Qwen3-4B DFlash proposer, \method{} is
1.483$\times$ and 1.246$\times$ faster than matched optimized source-trunk reuse on
MATH-500, with identical source/relay sequences and task accuracy on all
1,000 paired examples. The same interface-transplantation rule yields
1.250$\times$ and 1.081$\times$ held-out gains with a structurally distinct
EAGLE-3 chain proposer, without changing official task accuracy.
Unloading the eliminated source trunk saves 6.82--7.80~GiB of peak allocated
GPU memory across the measured backends and scales. Across 16 frozen
family--target--task cells, raw \method{} has a 1.114$\times$ geometric-mean
speedup; an acceptance-aware Amdahl model predicts all 16 point-estimate
directions with 0.34\% mean absolute relative error. A conservative profile
policy gives 1.135$\times$ descriptive geometric-mean speedup while retaining
source reuse for every measured slowdown or statistically unresolved cell.
The relay costs below 0.8\% of request time, so acceptance retention---rather
than relay compute---determines the net gain. This report specifies the
algorithm, correctness scope, measured evidence, novelty boundary, and
publication plan.
\end{abstract}

\callout{Core finding}{A 52--66M parameter translator can remove a measured
27--40\% repeated source-transformer component while preserving the target
verifier as the sole authority. DFlash agrees byte-for-byte with source reuse;
EAGLE-3 finite-precision differences are scored and audited rather than hidden.}

\section{Problem and research claim}

DFlash is an ICML 2026 block-diffusion speculative decoder that proposes a
token block in one parallel draft call and conditions that call on target
hidden features \cite{chen2026dflash}. Like EAGLE-3 and other feature-based
drafters, a released checkpoint is coupled to the feature space used during
its training \cite{li2025eagle3}. This creates a deployment choice for a new
compatible target model:

\begin{itemize}
  \item train and store a new target-specific proposer;
  \item run the original source model to reconstruct the old proposer's
        conditioning state; or
  \item translate the new verifier's already-computed state into the old
        interface.
\end{itemize}

\method{} realizes the third option. Its narrow, testable claim is:

\begin{quote}
\emph{A frozen feature-conditioned speculative proposer can be transplanted
across compatible target scales with a small per-target interface translator,
eliminating the repeated source transformer while retaining exact target
verification and delivering a positive end-to-end speed and memory gain.}
\end{quote}

This is not a claim that verifier features, KV reuse, or generic drafters are
new. DFlash and DFlare inject target features \cite{chen2026dflash,zhang2026dflare};
HyperDFlash aligns a newly trained DFlash with DeepSeek-V4's native
hyper-connection interface \cite{lin2026hyperdflash};
AngelSpec co-specializes newly trained proposer structures and adapts
verification effort to online workload \cite{liu2026angelspec};
RepSpec improves target-specific drafter training through inference-time-merged
re-parameterization \cite{huo2026repspec};
KVShot studies hidden/KV reuse \cite{liu2026kvshot}; PARD and OmniDraft study
portable drafters \cite{an2026pard,ramakrishnan2025omnidraft}; and
SD-squared steers pretrained drafters with verifier features
\cite{berdoz2026sd2}. The contribution is the specific \emph{interface
emulation plus trunk elimination} construction and its measured systems
consequences for an already-trained, unchanged proposer.
Recent verifier-skipping and sparse-verification work instead attacks target
verification, either through an explicitly lossy skip policy or specialized
long-context sparse-attention kernels \cite{luo2026verifierskipping,wang2026specsa}.
Those interventions are complementary; no such result is used to extrapolate
the dense-SDPA RelaySpec measurements.
SPEED-Bench further shows that speculative-decoding performance depends on
domain, context length, concurrency, and serving engine
\cite{abramovich2026speedbench}. Accordingly, four independent ranks are used
only to shorten experiment wall time; every reported rate is a batch-one
request-microaverage, not aggregate serving throughput.

\section{System decomposition}

Let $S$ be the source model used to train a released proposer $D$, and $T$
be a new target verifier. DFlash expects the conditioned state
\begin{equation}
c_t^S = \operatorname{RMSNorm}\!\left(
W_D [h_t^{S,s_1};\ldots;h_t^{S,s_m}]\right),
\end{equation}
where $W_D$ and the draft network are frozen. The pinned EAGLE-3 evaluator
likewise consumes a frozen projection of five source depths, but its
post-fusion state is raw rather than normalized. A matched cross-scale baseline
runs all layers of $S$ on the prefix and every proposed block, even though $T$
also processes that block for verification. The source calculation exists
only to maintain $D$'s old interface.

For the measured Qwen3-4B proposer, the source baseline divides into:

\begin{enumerate}
  \item DFlash block proposal;
  \item source-trunk block processing to refresh $c_t^S$;
  \item full target block verification;
  \item longest-prefix acceptance, KV-cache crop/commit, and loop control.
\end{enumerate}

Across the frozen workloads, source-trunk processing is 38--40\% of the
Qwen3-8B path and about 30\% of the Qwen3-14B path. This large measured
fraction, not a FLOP proxy, is the opportunity targeted by \method{}.

\begin{center}
\fbox{\begin{minipage}{0.93\linewidth}
\centering
\textbf{One speculative cycle}\par\smallskip
Exact target taps $\longrightarrow$ target-specific linear relay
$\longrightarrow$ frozen DFlash/EAGLE proposer
$\longrightarrow$ draft tokens
$\longrightarrow$ full-target verification and commit\par\smallskip
\textcolor{orange}{Source-reuse baseline only:} committed tokens
$\longrightarrow$ separate Qwen3-4B trunk $\longrightarrow$ old proposer
context. \method{} removes this orange critical-path branch; it does not skip
the verifier.
\end{minipage}}
\end{center}

\section{Method}

\subsection{Target feature relay}

Choose the $m$ target layers dictated by the released proposer interface and
concatenate verifier states produced for committed tokens,
\begin{equation}
z_t^T=[h_t^{T,t_1};\ldots;h_t^{T,t_m}].
\end{equation}
The shared translator is deliberately simple:
\begin{equation}
\widehat c_t = N_{\mathrm{out}}\!\left(
R_\phi(N_{\mathrm{in}}(z_t^T))\right),
\qquad R_\phi\in\mathbb{R}^{d_D\times m d_T}.
\end{equation}
For DFlash, $N_{\mathrm{in}}$ and $N_{\mathrm{out}}$ are the matched
normalizations of its consumed interface. For the pinned EAGLE-3 backend they
are identities, preserving the amplitude of its raw post-fusion context. This
is not a free family-specific embellishment: the matched 8B ablation changes
only normalization and improves relative interface loss from 0.413 to 0.266,
acceptance retention from 69.0\% to 82.6\%, and end-to-end speed from
1.035$\times$ to 1.237$\times$. The rule is then frozen and applied to 14B
without retuning.
There is no middle-layer predictor, recurrence, ensemble, or confidence gate.
One bias-free full-rank linear map is sufficient to test the core systems
hypothesis. A low-rank map cannot improve the dominant path materially because
the full-rank relay is already below 0.8\% of reference request time, while it
can reduce interface fidelity and therefore accepted tokens.

For Qwen3-8B, the map is $(5\times4096)\rightarrow2560$ (52,428,800
parameters); for Qwen3-14B it is $(5\times5120)\rightarrow2560$
(65,536,000 parameters). The FP32 checkpoints are 200 and 250~MiB.

\subsection{Inference algorithm}

At batch size one and greedy decoding:

\begin{enumerate}
  \item \textbf{Prefill.} Run $T$ on the prompt, create its exact KV cache,
        extract selected verifier taps, and compute $\widehat c$.
  \item \textbf{Draft.} Give $\widehat c$, the current token, and masked
        positions to the frozen DFlash proposer. It predicts a block of $k$
        draft tokens in one forward pass.
  \item \textbf{Verify.} Run the ordinary full target $T$ once on the proposed
        block using its exact committed KV cache.
  \item \textbf{Accept.} Commit the longest draft prefix that agrees with the
        target's greedy predictions, append the target correction token, and
        crop speculative KV entries beyond the committed prefix.
  \item \textbf{Relay.} Extract target taps for the newly committed positions,
        apply $R_\phi$, and use the result to condition the next draft cycle.
  \item Repeat until EOS or the preregistered output cap.
\end{enumerate}

The source embedding and LM head are retained because they are part of the
released proposer's token interface. The 36-layer Qwen3-4B transformer trunk
is deleted after loading. This reduces measured steady allocation by
6.82~GiB without changing relay throughput materially.

\subsection{Training objectives}

The target $T$, source $S$, and proposer $D$ are frozen. The preregistered
Part-2 objective predicts the exact proposer-consumed interface using
coefficient-free relative error:
\begin{equation}
\mathcal{L}_{\mathrm{rel}} = \frac{1}{N}\sum_t
\frac{\|\widehat c_t-c_t^S\|_2^2}
{\max(\|c_t^S\|_2^2,\epsilon_{\mathrm{dtype}})}.
\end{equation}
The denominator clamp is fixed by the accumulation dtype. If
$\rho=\|\widehat c\|/\|c^S\|$ and $\theta$ is the angle between the vectors,
then each term is $\rho^2+1-2\rho\cos\theta$. The objective therefore measures
both magnitude and direction without an unexplained mixing coefficient.

The completed Part-1 checkpoint historically used raw MSE plus a cosine term
weighted by $0.1$. We do not justify that inherited decimal post hoc. Before
the full evaluation, we refit both target scales with identical data, steps,
seed, optimizer, architecture, and taps, changing only the objective. Relative
MSE is faster at both scales: 1.5227$\times$ versus 1.5095$\times$ over source
reuse at 8B and 1.2653$\times$ versus 1.2514$\times$ at 14B. Direct
prompt-paired comparisons favor relative MSE by 1.0085$\times$
[1.0005,1.0174] and 1.0114$\times$ [1.0046,1.0177], respectively. A
preregistered conditional bracket would have tested cosine weights $0.03$ and
$0.3$ only if $0.1$ won; it did not, so no cosine coefficient enters the final
method. The historical detached-teacher KL refinement remains an ablation: its
1.0005$\times$ same-prompt gain has 95\% CI [0.9981,1.0028] and is unresolved.

Relay fitting uses a deterministic 4,096-example subset of the
\code{DigitalLearningGmbH/MATH-lighteval} training split at pinned revision
\code{0530c786...}. It supplies text positions for source/target interface
regression, not answer supervision: no language-model weights are updated. An
operator-preserving normalized-hash audit collapses whitespace while retaining
mathematical symbols and finds zero exact fit/evaluation overlap across all
1,250 evaluation records.

\section{Correctness scope}

The relay is approximate as a predictor of the old source conditioning state,
but it is not trusted to approve output tokens. At every cycle, the full target
verifies every proposed position; only a verifier-agreeing prefix and the
target correction token are committed. Therefore relay error affects
acceptance length and runtime, not the authority that selects committed
tokens. This is the standard separation underlying lossless speculative
decoding \cite{leviathan2023speculative,timor2025dsi,zhou2026hsd}.

Formally, let $P_j$ be the committed prefix before cycle $j$, let an arbitrary
provider---source trunk or relay---induce draft block $d_{1:k}$, and let
$v_i=\arg\max T(\cdot\mid P_j,d_{<i})$. The verifier commits
$d_{1:m}$ for the largest $m$ satisfying $d_i=v_i$ for every $i\leq m$, then
commits the target token $v_{m+1}$ when the cycle is not terminated. Hence the
new prefix is exactly the prefix selected by the target block-verification
path, regardless of how $d_{1:k}$ was produced. Induction on $j$ proves that
replacing the interface provider cannot give the relay token authority. The
implementation additionally crops target and proposer KV entries beyond the
committed prefix after every rejection; cache-conformance tests exercise that
invariant. This proposition is exact in arithmetic. It does not promise
bitwise equivalence between different target kernel shapes.

The implemented claim is greedy decoding. Within the same target block-verifier
path, source reuse and \method{} produced identical token sequences in every
paired run. A one-token autoregressive loop can differ at floating-point ties
because block and one-token calls select different kernel shapes; task
accuracy is therefore reported independently. Exact sampling would require
the standard distribution-correct rejection/correction rule and is not used
in these experiments.

\section{Experimental protocol}

\begin{table}[H]
\centering
\small
\begin{tabularx}{\linewidth}{@{}lX@{}}
\toprule
Item & Frozen choice \\
\midrule
Targets & Qwen3-8B and Qwen3-14B, pinned Hugging Face revisions \\
Proposer & \code{z-lab/Qwen3-4B-DFlash-b16}, pinned code/checkpoint \\
Reference & Native target AR; source-trunk reuse; target-specific DFlash-8B where released \\
Mode & Qwen3 non-thinking; greedy; batch one; BF16; SDPA \\
Tasks & MATH-500, GSM8K-128, HumanEval-164, EvalPlus MBPP-378, MT-Bench 80 conversations/160 turns \\
Output cap & 2,048 tokens for full task runs \\
Hardware & Exactly four NVIDIA RTX 6000 Ada 48~GB GPUs, one rank per GPU \\
Statistics & Prompt-paired 10,000-sample bootstrap 95\% intervals \\
Quality & Official Qwen math scorer and pinned EvalPlus base/Plus scorer; source/relay exact sequence match \\
Systems & End-to-end tokens/s, p50/p95 request latency, decode tokens/s, acceptance, per-position survival, target/draft calls, CUDA regions, peak/steady memory; DFlash TTFT measured, inherited EAGLE TTFT unavailable \\
\bottomrule
\end{tabularx}
\caption{Primary evaluation protocol. The task selection and 2,048-token cap
follow the ICML 2026 DFlash evaluation pattern \cite{chen2026dflash}.}
\end{table}

The four ranks evaluate disjoint prompts independently; no request is tensor
parallel across GPUs. After the excluded warm-up, each immutable manifest item
is executed once per paired method; uncertainty comes from whole-request
resampling, not repeated-prompt pseudo-replication. Tokens/s is total emitted tokens divided by summed
per-request latency, hence a micro-averaged single-request rate rather than
four-GPU aggregate serving throughput. Data-parallel sharding reduces
experimental wall time without multiplying the reported statistic. Within
each rank, method order rotates by request after a warmup so systematic thermal
or first-method effects do not favor the relay; speed intervals resample the
whole paired request and preserve its output-length coupling. The two
dependent turns of an MT-Bench conversation are resampled as one conversation
cluster. Official EvalPlus source/relay quality differences receive the same
task-paired bootstrap treatment rather than being compared as independent
binomial proportions.

The causal reference is deliberately optimized: after each verification cycle
it executes the Qwen3-4B source only for newly committed positions and only
through zero-based layer 33, the final consumed tap. Layers 34--35 and rejected
speculative positions have no causal path to proposer context and are omitted.
The naive full-source implementation appears only as an ablation, never as the
headline denominator. Native AR and released target-specific proposers are
separate deployment-frontier controls because they change the proposer rather
than only the interface provider.

Every fixed decision is registered before confirmatory scoring with one of six
provenance classes: formal derivation (F), peer-reviewed primary source (P),
pinned official artifact (O), disjoint-development selection (V), immutable
test observation (T), or administrative reproducibility value (A). Correctness
choices may only come from F; checkpoint interfaces from O; comparability
settings from P/O; and unresolved efficiency choices from V. Test observations
cannot retroactively choose their own method. Administrative budgets---for
example 4,096 fit examples and bootstrap seed 1729---are reported as controlled
budgets, not scientific optima. The released machine-readable protocol audits
every active configuration, rejects unregistered numeric constants, and
requires a rationale even for non-scientific numeric exemptions. This
hierarchy is itself the guard against circular backing: a weaker post-hoc test
result cannot override a stronger correctness or interface constraint.

MATH-500 is the primary task because it yields long non-thinking generations
and repeated verification cycles, exposing the source-trunk bottleneck without
the very long traces of explicit reasoning models. GSM8K tests shorter math,
HumanEval and MBPP test code, and MT-Bench tests long two-turn chat. The same
relay checkpoint is used across tasks at each target scale.

\section{Main results}

\subsection{Full MATH-500}

\begin{table}[H]
\centering
\small
\begin{tabular}{@{}lrrrrr@{}}
\toprule
Target & Native AR & Source reuse & \method{} & Relay/source & Accuracy S/R \\
\midrule
Qwen3-8B & 44.5 & 148.3 & \good{220.0} & \good{1.483$\times$ [1.475,1.492]} & 74.2/74.2 \\
Qwen3-14B & 26.7 & 110.4 & \good{137.6} & \good{1.246$\times$ [1.238,1.253]} & 79.6/79.6 \\
\bottomrule
\end{tabular}
\caption{End-to-end tokens/s and official accuracy. Source and relay sequences
match on 500/500 prompts at each scale.}
\end{table}

The released target-specific DFlash-8B reaches 241.2 tok/s on the matched
same-hardware control protocol; the selected relay reaches 91.2\% of that
cross-run contextual ceiling without training a new
proposer. The primary comparison is source reuse because it holds the frozen
proposer fixed and isolates the translated interface.
The 2,048-token cap is not hidden: 10.2\% of 8B MATH requests and 9.4\% of
14B requests hit it. Source and relay cap rates are identical because their
outputs are identical. Breadth cap rates are below 0.7\% in every completed
DFlash-8B task, so the cross-task speed result is not driven by systematic
early truncation.

\subsection{Second proposer family: DeepSpec EAGLE-3 chain}

The same interface-transplantation mechanism is evaluated with the pinned
DeepSpec EAGLE-3 \code{ttt7} evaluator. This is a length-7 autoregressive
proposal chain, not the dynamic-tree serving implementation, so every number is
compared only with its matched DeepSpec source-reuse path.

\begin{table}[H]
\centering
\small
\begin{tabular}{@{}lrrrrrr@{}}
\toprule
Target & Source & Relay & Relay/source & Source share & Retention & Accuracy S/R \\
\midrule
Qwen3-8B & 83.05 & 103.77 & \good{1.250$\times$ [1.244,1.256]} & 33.1\% & 84.0\% & 72.863/72.863 \\
Qwen3-14B & 64.93 & 70.20 & \good{1.081$\times$ [1.076,1.087]} & 27.1\% & 79.1\% & 79.487/79.487 \\
\bottomrule
\end{tabular}
\caption{Held-out 468-problem MATH complement. Throughput is end-to-end
tokens/s; intervals resample paired requests.}
\end{table}

The development-only Amdahl forecasts were 1.238$\times$ and 1.090$\times$;
held-out observations differ by only +0.9\% and -0.8\%. Byte agreement is
464/468 and 466/468. Every mismatch is retained in a per-problem audit and all
preserve the official correctness outcome, giving a paired accuracy delta and
interval of exactly zero at both scales.

\subsection{Deployment-frontier controls}

\begin{table}[H]
\centering
\scriptsize
\begin{tabular}{@{}llrrrrrr@{}}
\toprule
Family & Target & Native AR & Native prop. & Source & Relay & Relay/source & Relay/native \\
\midrule
DFlash & 8B & 44.52 & 241.25 & 148.33 & 219.99 & 1.483$\times$ & 91.19\% \\
DFlash & 14B & 26.71 & --- & 110.45 & 137.60 & 1.246$\times$ & --- \\
EAGLE-3 & 8B & 44.45 & 114.28 & 82.83 & 103.37 & 1.248$\times$ & 90.46\% \\
EAGLE-3 & 14B & 26.71 & 78.78 & 64.97 & 70.27 & 1.082$\times$ & 89.20\% \\
\bottomrule
\end{tabular}
\caption{Full MATH-500 end-to-end tok/s. Source/relay is the paired causal
comparison. Native controls run in separate memory-safe processes under the
same hardware/protocol and are contextual deployment-frontier comparisons.}
\end{table}

The released target-specific proposers remain faster because their interface
is native to the target. The portability result is that a 52--66M-parameter
relay recovers 89--91\% of that matched target-specific throughput where a
released control exists, while avoiding target-specific proposer retraining.
Source reuse, relay, and the corresponding block verifier have identical
official MATH correctness at each EAGLE scale; DFlash source and relay are
also identical. Native one-token AR accuracy is reported separately because
BF16 block and token kernel shapes can diverge at near-tied logits.

\subsection{Cross-task generalization}

\begin{table}[H]
\centering
\scriptsize
\resizebox{\textwidth}{!}{%
\begin{tabular}{@{}lllrrrrr@{}}
\toprule
Family & Target & Task & $N$/clusters & Source & Relay & Relay/source [95\% CI] & Exact \\
\midrule
EAGLE-3 & 8B & GSM8K & 128/128 & 85.99 & 108.45 & \good{1.261 [1.244,1.279]} & 99.22\% \\
EAGLE-3 & 8B & HumanEval & 164/164 & 71.75 & 81.10 & \good{1.130 [1.117,1.143]} & 96.95\% \\
EAGLE-3 & 8B & MBPP & 378/378 & 69.86 & 82.10 & \good{1.175 [1.166,1.185]} & 97.62\% \\
EAGLE-3 & 8B & MT-Bench & 160/80 & 42.27 & 49.45 & \good{1.170 [1.150,1.191]} & 89.38\% \\
EAGLE-3 & 14B & GSM8K & 128/128 & 67.81 & 73.41 & \good{1.083 [1.068,1.098]} & 100\% \\
EAGLE-3 & 14B & HumanEval & 164/164 & 55.38 & 50.58 & 0.913 [0.901,0.926] & 99.39\% \\
EAGLE-3 & 14B & MBPP & 378/378 & 54.55 & 50.49 & 0.926 [0.915,0.936] & 98.68\% \\
EAGLE-3 & 14B & MT-Bench & 160/80 & 33.83 & 34.57 & 1.022 [0.9999,1.045] & 96.88\% \\
\midrule
DFlash & 8B & GSM8K & 128/128 & 116.87 & 166.39 & \good{1.424 [1.404,1.444]} & 100\% \\
DFlash & 8B & HumanEval & 164/164 & 118.84 & 145.00 & \good{1.220 [1.202,1.238]} & 100\% \\
DFlash & 8B & MBPP & 378/378 & 117.77 & 151.55 & \good{1.287 [1.273,1.300]} & 100\% \\
DFlash & 8B & MT-Bench & 160/80 & 60.85 & 81.62 & \good{1.341 [1.318,1.364]} & 100\% \\
DFlash & 14B & GSM8K & 128/128 & 88.04 & 97.07 & \good{1.103 [1.084,1.121]} & 100\% \\
DFlash & 14B & HumanEval & 164/164 & 87.34 & 77.61 & 0.889 [0.874,0.904] & 100\% \\
DFlash & 14B & MBPP & 378/378 & 88.02 & 83.98 & 0.954 [0.942,0.966] & 100\% \\
DFlash & 14B & MT-Bench & 160/80 & 45.32 & 49.30 & \good{1.088 [1.065,1.108]} & 100\% \\
\bottomrule
\end{tabular}}
\caption{Frozen cross-task matrix. Throughput is end-to-end tok/s. Both
MT-Bench turns are resampled as one 80-conversation cluster; it has no judge
quality claim. Bold intervals lie strictly above one.}
\end{table}

Across all 16 cells, raw \method{} has a 1.1135$\times$ equal-cell geometric
mean and 11/16 strictly positive intervals. Official source/relay quality is
identical in every scored GSM8K, HumanEval, and MBPP cell. The EAGLE textual
differences are retained in mismatch audits; they do not change any paired
official score. The negative 14B code cells are part of the result rather than
being filtered after measurement.

\section{Amdahl analysis: where the speedup comes from}

Normalize optimized source-reuse latency to one. Let $p_C$ be its
cycle-dependent target-verification plus frozen-proposer share, $p_O$ the
non-source residual share, and $p_R$ the relay cost measured relative to that
reference latency. Let $a_S$ and $a_R$ denote mean committed tokens per cycle
for source reuse and RelaySpec. Removing the source trunk but requiring
$a_S/a_R$ as many cycles gives
\begin{equation}
\frac{L_R}{L_S}=\frac{a_S}{a_R}p_C+p_O+p_R,
\qquad
\frac{a_R}{a_S}>\frac{p_C}{1-p_O-p_R}
\end{equation}
as the strict no-slowdown condition. This is the selection rule used by the
development gate: acceptance must exceed the measured break-even, and the
paired-bootstrap lower bound of observed speed must also exceed one. It is not
a fitted threshold.

The deployable mechanism is therefore profile-gated rather than unconditional.
For each target--proposer--workload cell, a small registered calibration run
measures these terms and selects the relay provider only when both algebraic
break-even and the paired speed interval pass; otherwise the identical backend
keeps the source provider. This choice occurs outside token verification and
cannot alter target authority. All evaluated cells remain in the paper so the
gate is tested as a systems model rather than used as a reporting filter.

For the simpler equal-acceptance limit, let $p$ be the measured source-trunk
share and $r$ the measured replacement relay share. The component-profile
bound reduces to
\begin{equation}
S_{\mathrm{profile}}=\frac{1}{1-p+r}.
\end{equation}
The observed speed includes extra verifier cycles caused by lower acceptance;
the full model above accounts for that effect through $a_S/a_R$.

\begin{table}[H]
\centering
\scriptsize
\resizebox{\textwidth}{!}{%
\begin{tabular}{@{}lllrrrrrr@{}}
\toprule
Family & Target & Task & Source & Relay & Retention & Break-even & Predicted & Observed \\
\midrule
EAGLE-3 & 8B & GSM8K & 33.69\% & 0.38\% & 83.84\% & 65.88\% & 1.2644 & 1.2613 \\
EAGLE-3 & 8B & HumanEval & 33.46\% & 0.42\% & 75.36\% & 66.23\% & 1.1345 & 1.1304 \\
EAGLE-3 & 8B & MBPP & 33.56\% & 0.41\% & 78.28\% & 66.06\% & 1.1799 & 1.1753 \\
EAGLE-3 & 8B & MT-Bench & 33.20\% & 0.39\% & 78.37\% & 66.46\% & 1.1745 & 1.1698 \\
EAGLE-3 & 14B & GSM8K & 27.36\% & 0.41\% & 78.80\% & 72.31\% & 1.0872 & 1.0826 \\
EAGLE-3 & 14B & HumanEval & 27.23\% & 0.47\% & 66.62\% & 72.60\% & 0.9194 & 0.9132 \\
EAGLE-3 & 14B & MBPP & 27.26\% & 0.47\% & 67.44\% & 72.49\% & 0.9320 & 0.9256 \\
EAGLE-3 & 14B & MT-Bench & 26.83\% & 0.41\% & 75.08\% & 72.94\% & 1.0286 & 1.0221 \\
DFlash & 8B & GSM8K & 40.00\% & 0.56\% & 85.68\% & 59.60\% & 1.4228 & 1.4237 \\
DFlash & 8B & HumanEval & 39.46\% & 0.64\% & 73.96\% & 60.31\% & 1.2203 & 1.2201 \\
DFlash & 8B & MBPP & 39.57\% & 0.61\% & 77.84\% & 60.07\% & 1.2866 & 1.2868 \\
DFlash & 8B & MT-Bench & 38.52\% & 0.56\% & 82.73\% & 61.27\% & 1.3408 & 1.3412 \\
DFlash & 14B & GSM8K & 30.92\% & 0.59\% & 76.39\% & 68.85\% & 1.1063 & 1.1025 \\
DFlash & 14B & HumanEval & 30.46\% & 0.72\% & 61.88\% & 69.52\% & 0.8926 & 0.8886 \\
DFlash & 14B & MBPP & 30.59\% & 0.67\% & 66.10\% & 69.18\% & 0.9567 & 0.9541 \\
DFlash & 14B & MT-Bench & 29.80\% & 0.56\% & 76.69\% & 70.10\% & 1.0917 & 1.0877 \\
\bottomrule
\end{tabular}}
\caption{Measured cross-task Amdahl decomposition. Predictions match all
16 point-estimate directions with 0.34\% mean absolute relative error.}
\end{table}

This table supplies the central causal evidence. Merely making the relay
smaller cannot materially improve the system because it is already below
0.8\%. The useful optimization direction is better interface fidelity that
increases accepted tokens without restoring the source trunk. A conservative
workload-level policy requires both predicted speedup above one and a paired
speed lower bound above one. It keeps source reuse for both 14B code tasks and
the statistically unresolved EAGLE-14B MT-Bench cell, yielding a descriptive
1.1354$\times$ equal-cell geometric mean without selecting a measured
slowdown. This is calibrated routing, not a free per-prompt oracle.

\section{Ablations and resource evidence}

\subsection{Draft block size}

\begin{table}[H]
\centering
\small
\begin{tabular}{@{}rrrrr@{}}
\toprule
Block & Source tok/s & Relay tok/s & Relay/source & Target DFlash-8B \\
\midrule
8 & 103.7 & 164.9 & 1.590$\times$ & 175.1 \\
16 & \good{137.7} & \good{210.1} & 1.526$\times$ & \good{229.9} \\
32 & 88.7 & 131.4 & 1.480$\times$ & 207.7 \\
\bottomrule
\end{tabular}
\caption{Matched 64-prompt MATH ablation. Block 8 maximizes relative gain,
but the released b16 proposer's native block 16 maximizes actual throughput.}
\end{table}

All speculative methods match the source sequence on 64/64 prompts for every
block. The primary protocol keeps block 16 rather than selecting the largest
relative number after evaluation.

\subsection{Training objective}

\begin{table}[H]
\centering
\small
\begin{tabular}{@{}lrrr@{}}
\toprule
Target/objective & Relay/source & 95\% CI & Mean accepted \\
\midrule
8B relative MSE & \good{1.5227$\times$} & [1.4876,1.5581] & 6.961 \\
8B MSE + 0.1 cosine & 1.5095$\times$ & [1.4763,1.5445] & 6.910 \\
14B relative MSE & \good{1.2653$\times$} & [1.2426,1.2894] & 6.380 \\
14B MSE + 0.1 cosine & 1.2514$\times$ & [1.2287,1.2761] & 6.330 \\
\bottomrule
\end{tabular}
\caption{Matched 32-prompt objective selection before full evaluation. Every
candidate passes the executable gate and matches 32/32 source sequences;
coefficient-free relative MSE is selected at both scales.}
\end{table}

\subsection{Memory residency and adaptation cost}

Deleting the source trunk after retaining its embedding/head changes DFlash-8B
allocated memory before requests from 23.92 to 17.10~GiB and peak allocation
from 24.16 to 17.34~GiB: 6.82~GiB saved (28.5\% steady, 28.2\% peak). Relay
throughput changes from 210.22 to 209.66 tok/s (99.74\% retained). Independent
single-method EAGLE-3 processes reproduce the mechanism: the 8B relay saves
7.45~GiB before requests and 7.80~GiB at peak (30.35\%/30.51\%), while the 14B
relay saves 7.42~GiB and 7.63~GiB (20.17\%/20.30\%). These isolated
micro-throughputs are contextual checks, not substituted for paired speed
estimates.

The 8B feature relay trains for 1,024 distributed steps over 4,096 math
examples in 85.7 seconds of instrumented time (107 seconds including setup) on
four Ada GPUs. Conservative proposal refinement takes 101.1 seconds and is
retained only as the ablation above. This is
the adaptation-efficiency result: a 200~MiB per-target adapter and minutes of
training replace a multi-GiB source trunk or a separately trained proposer.

\section{Relation to current literature and novelty audit}

\begin{longtable}{@{}p{0.18\linewidth}p{0.37\linewidth}p{0.37\linewidth}@{}}
\toprule
Work & What it already establishes & Distinction of \method{} \\
\midrule
\endhead
DFlash \cite{chen2026dflash} & Target-feature-conditioned block diffusion and strong lossless speedups. & Reuses its frozen conditioning contract for other target scales; does not redesign DFlash. \\
DFlare \cite{zhang2026dflare} & Per-draft-layer target fusion, deeper drafts, 2.4M-sample target-specific training. & One small external translator; no draft-layer change; minutes of per-target fitting. \\
HyperDFlash \cite{lin2026hyperdflash} & Target-native feature-boundary alignment and an inherited gated reducer for a newly trained DeepSeek-V4 drafter. & Keeps the old proposer frozen, translates a different target into its old interface, and removes a separately executed source trunk; reports component and memory effects. \\
AngelSpec \cite{liu2026angelspec} & Co-specialized MTP/DFly proposers and cost-aware online verification allocation under serving load. & Reuses an unchanged proposer across targets; current claims are batch-one latency and are not extrapolated to serving throughput. \\
RepSpec \cite{huo2026repspec} & Training-time structural re-parameterization that is merged into a target-specific EAGLE drafter for inference. & Trains only the external cross-target interface and leaves the old proposer weights and structure unchanged. \\
SPEED-Bench \cite{abramovich2026speedbench} & Standardized domain-diverse and concurrency-aware speculative-decoding evaluation on production engines. & Current evidence is deliberately batch-one latency; production throughput is outside the claim until RelaySpec is integrated into such an engine. \\
EAGLE-3 \cite{li2025eagle3} & Strong target-trained hidden-feature drafter and evaluation template. & Transplants an existing proposer rather than training a new target drafter. \\
KVShot \cite{liu2026kvshot} & Hidden/KV reuse diagnosis and long-range acceptance limits. & Removes a measured source-transformer critical-path component and reports end-to-end/memory gains. \\
PARD / OmniDraft \cite{an2026pard,ramakrishnan2025omnidraft} & Portable or target-independent drafters. & Per-target interface adapter for a released target-specific proposer, not a newly designed universal drafter. \\
SD-squared \cite{berdoz2026sd2} & Verifier-feature steering of pretrained autoregressive drafters. & Emulates a proposer's learned conditioning interface and replaces a whole repeated trunk. \\
C2C-style transfer & Cross-model latent/KV communication for multi-model reasoning. & Same-request speculative drafting; speedup comes from verifier-cycle trunk elimination, not semantic agent communication. \\
Cross-model KV transfer \cite{heo2026crossmodelkv} & Closed-form source-to-receiver KV mapping skips receiver prefill, with architecture constraints and measured accuracy retention. & Target-to-proposer interface mapping removes a separate decode-time source trunk; it does not synthesize verifier KV, and full-target verification retains token authority. \\
CaDDTree / DARTree \cite{zhang2026caddtree,li2026dartree} & Candidate-tree and cost-aware verification budget selection. & Orthogonal proposer-portability layer; current evaluation keeps a single block trajectory. \\
Speculative KV coding \cite{finn2026speculativekvcoding} & Non-peer-reviewed technical note on entropy coding an exact KV cache. & No cache compression or arithmetic coding; predicts only the frozen proposer's conditioning state. \\
\bottomrule
\end{longtable}

The novelty claim survives only in this precise form. ``Using target hidden
states,'' ``mapping between models,'' ``portable drafting,'' and
``acceptance-aware training'' are all too broad and conflict with prior work
\cite{zou2026vsd,an2026pard2,lin2026eda}. The paper should lead with measured
trunk elimination and frozen-proposer transplantation.

\section{ICLR paper package}

\subsection{Recommended paper story}

\begin{enumerate}
  \item \textbf{Observation:} cross-scale reuse of a feature-conditioned
        proposer introduces a redundant source transformer that occupies
        27--40\% of measured request time.
  \item \textbf{Method:} emulate the frozen proposer's interface with one
        target-specific linear relay fed by already-computed verifier taps.
  \item \textbf{Guarantee:} the relay proposes but never verifies; the full
        target remains the only commit authority.
  \item \textbf{Result:} clean DFlash gains of
        1.483$\times$/1.246$\times$ and independent EAGLE-3 gains of
        1.250$\times$/1.081$\times$ on MATH, with unchanged paired task
        accuracy; breadth and isolated memory evidence complete the matrix.
  \item \textbf{Explanation:} Amdahl decomposition predicts the ceiling and
        acceptance retention explains task/scale variation.
\end{enumerate}

\subsection{Core figures and tables}

The submission should contain: (1) a cycle diagram contrasting source reuse
and relay; (2) the two-scale main table; (3) the Amdahl stacked runtime
breakdown; (4) speedup versus acceptance-retention scatter with the profile
bound; (5) block-size and training-objective ablations; and (6) adapter
size/training-time/memory table. Per-position acceptance survival belongs in
the appendix, following alignment analyses such as GRIFFIN
\cite{hu2025griffin}.

\subsection{Evidence already complete}

\begin{itemize}
  \item two target scales, five workloads, and official MATH/EvalPlus scoring;
  \item native AR, matched source reuse, and released target-specific DFlash-8B;
  \item paired CIs, sequence identity, cap rates, acceptance, calls,
        component timings, source snapshots, pinned revisions, and GPU logs;
  \item memory loaded/unloaded measurement, block-size sweep, feature versus
        proposal-aligned objective, and exactly-four-GPU allocation checks.
\end{itemize}

\subsection{Second-family evidence}

The highest-value mechanism extension is complete: the same transplantation
rule works with the pinned DeepSpec EAGLE-3 chain proposer at both 8B and 14B,
with zero official MATH accuracy delta and held-out gains of 1.250$\times$ and
1.081$\times$. This separates the paper's claim from DFlash-specific diffusion
details. A production SGLang or vLLM concurrency study remains future serving
work; it is not needed to establish the batch-one causal mechanism and is not
silently mixed with the current PyTorch results.

\section{Reproducibility and artifact map}

Every GPU generation job records the exact config, immutable manifest and
hash, source snapshot and SHA-256 list, model/code revisions, Python/CUDA/
PyTorch/Transformers versions, device topology, rank-separated JSONL rows, and
10-second GPU telemetry. The principal local artifacts are:

\begin{itemize}
  \item \path{reports/final/dflash-8b-math500} and
        \path{reports/final/dflash-14b-math500};
  \item \path{reports/final/eagle3-8b-math500} and
        \path{reports/final/eagle3-14b-math500};
  \item all four frozen breadth directories plus the two corrected EAGLE
        two-turn MT-Bench directories under \path{reports/final};
  \item \path{reports/final/BREADTH_MATRIX.json} and its generated Markdown;
  \item \path{reports/final/EAGLE3_8B_MEMORY.md} and
        \path{reports/final/EAGLE3_14B_MEMORY.md};
  \item final ledger and status: \path{reports/FINAL_RESULTS.md} and
        \path{reports/EXECUTION_STATUS.md};
  \item mechanism and cost reports: \path{reports/amdahl-analysis.md} and
        \path{reports/adapter-cost.md}.
\end{itemize}

The complete local suite and final remote suite pass 118 tests. The remote
allocation exposes exactly four RTX 6000 Ada GPUs; its allocation
and pytest proof are stored under
\path{reports/final/remote-4gpu-verification}. GPU launchers fail fast unless
the world size and visible CUDA device count are exactly four.

\section{Conclusion}

\method{} is a simple systems method with a clear bottleneck, a compact
algorithm, and measured positive results. It does not attempt to approximate
the expensive target verifier. Instead, it removes a different large cost:
the source transformer executed only to satisfy an old proposer's conditioning
contract. A tiny translator turns verifier work that is already mandatory
into the next draft context. The main evidence---paired end-to-end speed,
official quality, audited finite-precision agreement, memory savings, and
Amdahl accounting---supports an ICLR paper centered on frozen-proposer
portability and critical-path elimination. Raw relay use is not universally
faster at 14B; the measured cost model predicts that boundary and a
conservative calibrated provider keeps the source path when positive speed
evidence is absent.

\bibliographystyle{plain}
\bibliography{references}

\end{document}
